\section{基于有序Logistic回归的声环境满意度影响因素分析}

\subsection{变量定义}

\subsubsection{自变量卡方检验}

有序Logistic回归对自变量的筛选需兼顾统计显著性与实际意义。本研究选取性别、年龄、文化程度、所在城市、居住区域、住宅隔音配置、居家办公共7项分类变量，分别与居民声环境满意度开展Pearson卡方独立性检验，以 $P<0.05$ 作为变量入选标准，依托467份有效问卷的样本频数构建交叉列联表进行筛选。

\begin{table}[H]
\centering
\caption{协变量筛选表}
\rowcolors{2}{tablight}{white}
\begin{tabular}{M{5cm}M{4cm}M{4cm}}
\bluetoprule
\rowcolor{tabhead}
\thead{变量} & \thead{$P$ 值} & \thead{是否保留} \\
\hline
性别 & 0.392 & 剔除 \\
年龄 & $<0.05$ & 保留 \\
文化程度 & $<0.05$ & 保留 \\
所在城市 & $<0.05$ & 保留 \\
居住区域 & $<0.05$ & 保留 \\
住宅隔音配置 & $<0.05$ & 保留 \\
居家办公 & 0.287 & 剔除 \\
\bluebottomrule
\end{tabular}
\end{table}

\subsubsection{变量设计}

被解释变量为居民声环境满意度，采用五级李克特量表有序赋值：1=非常不满意，2=不太满意，3=一般，4=比较满意，5=非常满意。解释变量涵盖人口特征、居住环境与噪声感知三大类，其中分类变量统一设置为虚拟变量并选定参照组，连续变量采用公因子得分原值纳入。各变量具体说明如表\ref{tab:vars}所示。

\begin{table}[H]
\centering
\caption{变量设计表}
\label{tab:vars}
\resetrowcolors
\begin{tabular}{M{1.8cm}M{3.4cm}M{2.0cm}p{7.0cm}}
\bluetoprule
\rowcolor{tabhead}
\thead{变量分类} & \thead{变量名称} & \thead{变量类型} & \multicolumn{1}{c}{\thead{变量赋值与说明}} \\
\hline
被解释变量 & Y：声环境满意度 & 有序分类 & 1=非常不满意；2=不太满意；3=一般；4=比较满意；5=非常满意 \\
\hline
\multirow{2}{*}{人口特征} & 年龄（Age） & 虚拟变量 & 以18-30岁为参照组，设4个虚拟变量：$Age_1$=18岁及以下，$Age_2$=31-45岁，$Age_3$=46-60岁，$Age_4$=61岁及以上 \\ \cline{2-4}
 & 文化程度（Edu） & 虚拟变量 & 以本科及以上为参照组，设4个虚拟变量：$Edu_1$=小学及以下；$Edu_2$=初中；$Edu_3$=高中/中专/技校；$Edu_4$=大专 \\
\hline
\multirow{3}{*}{居住环境} & 所在城市（City） & 虚拟变量 & 以温州为参照组，设2个虚拟变量：$City\_HZ$=杭州；$City\_NB$=宁波 \\ \cline{2-4}
 & 居住区域（Area） & 虚拟变量 & 以乡镇、农村为参照组，设3个虚拟变量：$Area_1$=城市主城区；$Area_2$=城市郊区；$Area_3$=县城 \\ \cline{2-4}
 & 隔音墙配置（Wall） & 二分类 & 1=住宅配备隔音墙；0=无隔音墙或不清楚 \\
\hline
\multirow{4}{*}{噪声感知} & $F_1$：身心健康与认知干扰 & 连续变量 & 因子综合得分，涵盖噪声对入睡、注意力、精神状态及听觉健康的影响 \\ \cline{2-4}
 & $F_2$：情绪与社会关系影响 & 连续变量 & 因子综合得分，涵盖噪声对睡眠中断、人际关系、身体健康及极端情绪的影响 \\ \cline{2-4}
 & $F_3$：典型环境噪声源干扰 & 连续变量 & 因子综合得分，涵盖交通噪声、施工噪声及公共活动噪声的干扰程度 \\ \cline{2-4}
 & $F_4$：生活场景噪声源干扰 & 连续变量 & 因子综合得分，涵盖周围人生活噪声、商业经营噪声及动物噪声的干扰程度 \\
\bluebottomrule
\end{tabular}
\end{table}

\subsection{有序Logistic回归模型设定}

被解释变量为有序分级数据，等级之间不存在等距关系，不满足普通最小二乘线性回归的数据前提；同时五级满意度不宜简单合并为二分类变量造成信息损耗，因此选用比例优势假设下的累积有序Logistic模型，设 $Y$ 为居民声环境满意度（$Y=1,2,3,4,5$ 分别对应“非常不满意”至“非常满意”），$X_i$ 为第 $i$ 个解释变量，累积概率 $P(Y\le j)$ 表示满意度不超过等级 $j$ 的概率。对于包含 $p$ 个解释变量的模型，第 $j$ 个累积Logit方程的形式为：

\begin{equation}
\mathrm{Logit}\bigl[P(Y\le j)\bigr]=\alpha_j-\sum_{i=1}^{p}\beta_i X_i,\qquad j=1,2,3,4
\end{equation}

式中：$\alpha_j$ 为第 $j$ 个累积切分点的截距参数（阈值），满足 $\alpha_1<\alpha_2<\alpha_3<\alpha_4$；$\beta_i$ 为第 $i$ 个解释变量 $X_i$ 的回归系数，反映在其他变量保持不变时，$X_i$ 每变动一个单位所引起的累积Logit值的平均变化量。$\beta_i>0$ 表示 $X_i$ 取值越大，居民倾向于选择更高满意度等级的概率越大（正向影响）；$\beta_i<0$ 则表示 $X_i$ 取值越大，居民倾向于选择更低满意度等级的概率越大（负向影响）。

将本研究选取的全部解释变量代入上述模型，展开后的完整形式为：

\begin{equation}
\begin{aligned}
\mathrm{Logit}\bigl[P(Y\le j)\bigr]={}&\alpha_j-\beta_1 Age_1-\beta_2 Age_2-\beta_3 Age_3-\beta_4 Age_4\\
&-\beta_5 Edu_1-\beta_6 Edu_2-\beta_7 Edu_3-\beta_8 Edu_4\\
&-\beta_9 City\_HZ-\beta_{10} City\_NB\\
&-\beta_{11} Area_1-\beta_{12} Area_2-\beta_{13} Area_3\\
&-\beta_{14} Wall\\
&-\beta_{15} F_1-\beta_{16} F_2-\beta_{17} F_3-\beta_{18} F_4,\qquad j=1,2,3,4
\end{aligned}
\end{equation}

模型共包含4个累积Logit方程，对应4个切分点。每个方程中18个解释变量的回归系数 $\beta_1$ 至 $\beta_{18}$ 完全相同，仅截距项 $\alpha_j$ 随切分点变化——这是比例优势假设的核心约束：解释变量对满意度从“非常不满意”过渡到“不太满意”的效应，与从“比较满意”过渡到“非常满意”的效应，在统计意义上等价。

由累积概率可进一步反推各满意度等级的具体概率：

\begin{equation}
\left\{\begin{aligned}
P(Y=1)&=P(Y\le1)\\
P(Y=2)&=P(Y\le2)-P(Y\le1)\\
P(Y=3)&=P(Y\le3)-P(Y\le2)\\
P(Y=4)&=P(Y\le4)-P(Y\le3)\\
P(Y=5)&=1-P(Y\le4)
\end{aligned}\right.
\end{equation}

模型参数求解采用最大似然估计法，通过迭代算法寻找使样本观测值出现概率最大的参数组合。本研究利用SPSS完成全部参数估计与模型检验。

\subsection{模型检验}

为保证回归结果的稳健性与可靠性，本研究依次开展平行性检验、多重共线性检验与模型拟合优度检验，三项检验从模型前提假设、变量间关系与整体解释力三个层面为回归分析的有效性提供保障。

\subsubsection{平行性检验}

平行性检验（Test of Parallel Lines），亦称比例优势假设检验，是有序Logistic回归最为核心的前提条件。该假设要求：对于被解释变量的任意两个不同切分点，所有解释变量的回归系数均相等，即解释变量对累积Logit的影响不随切分点的变化而变化。若平行性假设不成立，则需改用广义有序Logit模型或多分类Logit模型。

本研究采用似然比检验法，比较约束模型（平行性假设成立）与非约束模型（广义模型，各切分点系数自由估计）的$-2$倍对数似然值差异。检验结果如表\ref{tab:parallel}所示。

\begin{table}[H]
\centering
\caption{似然比检验结果}
\label{tab:parallel}
\rowcolors{2}{tablight}{white}
\begin{tabular}{M{4.8cm}M{2.6cm}M{2.0cm}M{1.8cm}M{2.0cm}}
\bluetoprule
\rowcolor{tabhead}
\thead{模型} & \thead{$-2$ Log Likelihood} & \thead{$\chi^2$} & \thead{自由度} & \thead{显著性} \\
\hline
原假设模型（比例优势成立） & 1227.952 & — & — & — \\
广义模型（各切分点系数自由） & 1140.812 & — & — & — \\
差值（广义 $-$ 原假设） & — & 87.14 & 54 & 0.003 \\
\bluebottomrule
\end{tabular}
\end{table}

检验结果显示，广义模型相对于原假设模型的似然比检验显著性 $P=0.003<0.05$，比例优势假设在严格检验下被拒绝，说明部分自变量的效应强度可能随满意度切分点而变化。考虑到该检验在样本量较大时对微小偏离较为敏感，本文在保留比例优势模型作为主模型的同时，进一步拟合广义有序模型进行敏感性分析，以判断这一偏离是否影响核心结论。

为进一步检验核心结论的稳健性，本文拟合广义有序模型（各切分点系数自由估计），四个噪声感知因子在各切分点上的系数如表\ref{tab:sensitivity}所示。

\begin{table}[H]
\centering
\caption{广义有序模型的各切分点系数（敏感性分析）}
\label{tab:sensitivity}
\rowcolors{2}{tablight}{white}
\begin{tabular}{M{2.6cm}M{2.6cm}M{2.6cm}M{2.6cm}M{2.6cm}}
\bluetoprule
\rowcolor{tabhead}
\thead{切分点} & \thead{$F_1$ 身心} & \thead{$F_2$ 情绪} & \thead{$F_3$ 环境源} & \thead{$F_4$ 生活源} \\
\hline
$Y\le1$ & $-0.337$ & $-0.152$ & $-0.631$ & $-0.433$ \\
$Y\le2$ & $-0.124$ & $-0.417$ & $-0.304$ & $-0.130$ \\
$Y\le3$ & $-0.119$ & $-0.625$ & $-0.458$ & $-0.355$ \\
$Y\le4$ & $-0.917$ & $-1.088$ & $-0.599$ & $-0.463$ \\
\bluebottomrule
\end{tabular}
\end{table}

由表\ref{tab:sensitivity}可见，$F_1\sim F_4$ 四个噪声感知因子在四个切分点上的系数全部为负、方向完全一致，表明“噪声影响越强、满意度越低”的核心结论在不同切分点下均成立，比例优势假设的偏离仅体现为效应强度的变化，而非影响方向的反转，核心结论稳健。鉴于比例优势模型在系数解释上的直观性，本文仍以其为主要模型，并以广义有序模型作为敏感性分析予以佐证。

\subsubsection{多重共线性检验}

多重共线性是指解释变量之间存在高度线性相关关系，将导致回归系数估计的标准误膨胀、参数显著性检验失效，甚至出现回归系数符号与理论预期相悖的问题。本研究采用方差膨胀因子（Variance Inflation Factor, VIF）对全部解释变量进行共线性诊断。VIF值越大，表明该变量与其他变量的共线性越严重；一般以 $VIF<10$ 为可接受的判断标准，更严格的标准为 $VIF<5$。

将全部解释变量构建普通线性回归模型，以容忍度（Tolerance）和VIF作为诊断指标，检验结果如表\ref{tab:vif}所示。

\begin{table}[H]
\centering
\caption{多重共线性检验结果}
\label{tab:vif}
\rowcolors{2}{tablight}{white}
\begin{tabular}{M{6.4cm}M{3.4cm}M{2.4cm}}
\bluetoprule
\rowcolor{tabhead}
\thead{变量} & \thead{容忍度（Tolerance）} & \thead{VIF} \\
\hline
$Age_1$（18 岁以下） & 0.779 & 1.283 \\
$Age_2$（31-45 岁） & 0.620 & 1.612 \\
$Age_3$（46-60 岁） & 0.656 & 1.525 \\
$Age_4$（61 岁及以上） & 0.721 & 1.387 \\
$Edu_1$（小学及以下） & 0.779 & 1.283 \\
$Edu_2$（初中） & 0.728 & 1.374 \\
$Edu_3$（高中 / 中专 / 技校） & 0.758 & 1.319 \\
$Edu_4$（大专） & 0.766 & 1.306 \\
$City\_HZ$（杭州） & 0.689 & 1.452 \\
$City\_NB$（宁波） & 0.645 & 1.550 \\
$Area_1$（城市主城区） & 0.542 & 1.846 \\
$Area_2$（城市郊区） & 0.559 & 1.788 \\
$Area_3$（县城） & 0.634 & 1.578 \\
$Wall$（隔音墙配置） & 0.920 & 1.087 \\
$F_1$（身心健康与认知干扰） & 0.946 & 1.057 \\
$F_2$（情绪与社会关系影响） & 0.978 & 1.022 \\
$F_3$（典型环境噪声源干扰） & 0.959 & 1.043 \\
$F_4$（生活场景噪声源干扰） & 0.884 & 1.131 \\
\bluebottomrule
\end{tabular}
\end{table}

检验结果显示，全部自变量的容忍度均大于0.5，VIF值介于1.022至1.846之间，远低于10的常用阈值，也均未超过更严格的标准5。其中居住区域变量（$Area_1$、$Area_2$、$Area_3$）的VIF值相对偏高，但仍在安全范围内。总体而言，变量间不存在严重的多重共线性问题，回归系数的估计具有无偏性与稳定性。

\subsubsection{模型拟合优度检验}

模型拟合优度用于评估回归模型整体对数据的解释能力。本研究从模型整体显著性与伪决定系数两个方面进行评估，检验结果如表\ref{tab:gof}所示。

\begin{table}[H]
\centering
\caption{模型拟合度检验结果}
\label{tab:gof}
\rowcolors{2}{white}{white}
\begin{tabular}{M{6.0cm}M{3.0cm}M{2.2cm}M{2.2cm}}
\bluetoprule
\rowcolor{tabhead}
\thead{检验指标} & \thead{数值} & \thead{自由度} & \thead{显著性} \\
\hline
\rowcolor{tablight}
\multicolumn{4}{l}{\textbf{模型拟合信息（Model Fitting Information）}} \\
仅截距模型 $-2$ Log Likelihood & 1349.544 & — & — \\
最终模型 $-2$ Log Likelihood & 1227.952 & — & — \\
似然比 $\chi^2$ & 121.592 & 18 & $<0.001$ \\
\rowcolor{tablight}
\multicolumn{4}{l}{\textbf{伪决定系数（Pseudo $R^2$）}} \\
Cox \& Snell $R^2$ & 0.229 & — & — \\
Nagelkerke $R^2$ & 0.243 & — & — \\
McFadden $R^2$ & 0.090 & — & — \\
\bluebottomrule
\end{tabular}
\end{table}

由结果可知：

\textbf{（1）模型整体显著性}：仅含截距的零模型与纳入全部解释变量的最终模型之间的似然比卡方值为121.592（df=18，$P<0.001$），拒绝“所有解释变量回归系数均为零”的原假设，表明至少有一个解释变量对声环境满意度具有显著影响，模型整体具有统计学意义。

\textbf{（2）伪决定系数}：Nagelkerke $R^2$ 为0.243，Cox \& Snell $R^2$ 为0.229，McFadden $R^2$ 为0.090。在社会科学的问卷调查研究领域，Nagelkerke $R^2$ 超过0.2表明模型具有一定的解释能力，但相对偏低，说明声环境满意度除受本文纳入变量影响外，还受到其他未观测因素的制约。

综合以上检验结果，本研究所构建的有序Logistic回归模型变量间不存在严重共线性，模型整体拟合效果可以接受。尽管比例优势假设在严格检验下被拒绝、伪决定系数偏低，但敏感性分析显示核心结论稳健，模型所揭示的变量间关系仍具有统计意义，可用于识别影响声环境满意度的关键因素。

\subsection{回归结果分析}

\subsubsection{参数估计结果}

有序Logistic回归的参数估计结果按变量类别分列于表\ref{tab:params1}与表\ref{tab:params2}。前者给出模型切分点（阈值）与人口特征、居住环境变量的估计结果，后者单独给出四个噪声感知因子的估计结果。

\begin{table}[H]
\centering
\caption{模型切分点与人口特征、居住环境变量的参数估计结果}
\label{tab:params1}
\resizebox{\textwidth}{!}{%
\begin{tabular}{M{3.2cm}M{1.5cm}M{1.5cm}M{1.6cm}M{0.9cm}M{1.4cm}M{1.7cm}M{2.6cm}}
\bluetoprule
\rowcolor{tabhead}
\thead{变量} & \thead{估计值\ ($\beta$)} & \thead{标准误\ (SE)} & \thead{Wald\ $\chi^2$} & \thead{df} & \thead{$P$值} & \thead{OR值\ Exp($\beta$)} & \thead{OR的95\%置信区间} \\
\hline
\rowcolor{tablight}
\multicolumn{8}{l}{\textbf{阈值}} \\
$Y=1$（非常不满意） & -3.196 & 0.331 & 93.523 & 1 & $<0.001$ & — & — \\
$Y=2$（不太满意） & -0.865 & 0.292 & 8.779 & 1 & 0.003 & — & — \\
$Y=3$（一般） & 0.713 & 0.296 & 5.824 & 1 & 0.016 & — & — \\
$Y=4$（比较满意） & 2.433 & 0.331 & 53.873 & 1 & $<0.001$ & — & — \\
\rowcolor{tablight}
\multicolumn{8}{l}{\textbf{人口特征}} \\
$Age_1$：18 岁及以下 & -0.995 & 0.371 & 7.182 & 1 & 0.007 & 0.370 & 0.179-0.765 \\
$Age_2$：31-45 岁 & -0.321 & 0.227 & 1.998 & 1 & 0.158 & 0.725 & 0.465-1.132 \\
$Age_3$：46-60 岁 & -0.648 & 0.252 & 6.578 & 1 & 0.010 & 0.523 & 0.319-0.858 \\
$Age_4$：61 岁及以上 & -0.557 & 0.330 & 2.851 & 1 & 0.091 & 0.573 & 0.300-1.094 \\
$Edu_1$：小学及以下 & -0.082 & 0.351 & 0.055 & 1 & 0.815 & 0.921 & 0.463-1.834 \\
$Edu_2$：初中 & -0.521 & 0.286 & 3.314 & 1 & 0.069 & 0.594 & 0.339-1.041 \\
$Edu_3$：高中 / 中专 / 技校 & -0.477 & 0.250 & 3.642 & 1 & 0.056 & 0.620 & 0.380-1.013 \\
$Edu_4$：大专 & -0.648 & 0.260 & 6.209 & 1 & 0.013 & 0.523 & 0.314-0.871 \\
\rowcolor{tablight}
\multicolumn{8}{l}{\textbf{居住环境}} \\
$City\_HZ$：杭州 & -0.162 & 0.225 & 0.522 & 1 & 0.470 & 0.850 & 0.547-1.320 \\
$City\_NB$：宁波 & 0.082 & 0.216 & 0.144 & 1 & 0.704 & 1.085 & 0.711-1.656 \\
$Area_1$：城市主城区 & 0.191 & 0.261 & 0.535 & 1 & 0.464 & 1.211 & 0.725-2.021 \\
$Area_2$：城市郊区 & -0.217 & 0.255 & 0.728 & 1 & 0.394 & 0.805 & 0.488-1.326 \\
$Area_3$：县城 & -0.400 & 0.280 & 2.040 & 1 & 0.153 & 0.670 & 0.387-1.160 \\
$Wall$：配备隔音墙 & 0.670 & 0.201 & 11.140 & 1 & $<0.001$ & 1.953 & 1.318-2.894 \\
\bluebottomrule
\end{tabular}}
\end{table}

\begin{table}[H]
\centering
\caption{噪声感知因子的参数估计结果}
\label{tab:params2}
\resizebox{\textwidth}{!}{%
\begin{tabular}{M{3.2cm}M{1.5cm}M{1.5cm}M{1.6cm}M{0.9cm}M{1.4cm}M{1.7cm}M{2.6cm}}
\bluetoprule
\rowcolor{tabhead}
\thead{变量} & \thead{估计值\ ($\beta$)} & \thead{标准误\ (SE)} & \thead{Wald\ $\chi^2$} & \thead{df} & \thead{$P$值} & \thead{OR值\ Exp($\beta$)} & \thead{OR的95\%置信区间} \\
\hline
\rowcolor{tablight}
\multicolumn{8}{l}{\textbf{噪声感知因子}} \\
$F_1$：身心健康与认知干扰 & -0.256 & 0.092 & 7.841 & 1 & 0.005 & 0.774 & 0.647-0.926 \\
$F_2$：情绪与社会关系影响 & -0.583 & 0.095 & 37.851 & 1 & $<0.001$ & 0.558 & 0.464-0.672 \\
$F_3$：典型环境噪声源干扰 & -0.440 & 0.092 & 22.719 & 1 & $<0.001$ & 0.644 & 0.538-0.772 \\
$F_4$：生活场景噪声源干扰 & -0.309 & 0.094 & 10.731 & 1 & 0.001 & 0.734 & 0.610-0.883 \\
\bluebottomrule
\end{tabular}}
\end{table}

注：参照组：年龄（18-30岁）、文化程度（本科及以上）、城市（温州）、居住区域（乡镇/农村）。$OR$ 值 $<1$ 表示该变量取值越高，满意度越倾向于低等级（负向影响）；$OR$ 值 $>1$ 表示该变量取值越高，满意度越倾向于高等级（正向影响）。

\subsubsection{结果解读}

\textbf{（一）人口特征变量的影响}

年龄维度中，\textbf{18岁及以下组}的负向效应最强（$\beta=-0.995$，$P=0.007$），其满意度处于更高等级的累积优势仅为18-30岁参照组的\textbf{37.0\%}；\textbf{46-60岁组}次之（$\beta=-0.648$，$P=0.010$，累积优势\textbf{52.3\%}）。31-45岁组与61岁及以上组未通过显著性检验。

文化程度维度中，仅\textbf{大专组}达到显著水平（$\beta=-0.648$，$P=0.013$，累积优势\textbf{52.3\%}），其余三组均不显著。

综合来看，人口特征变量的\textbf{独立解释力有限}，显著项仅集中于青少年、46-60岁中年与大专学历三类群体，满意度更多取决于居民所处的声环境条件与主观噪声感知。

\textbf{（二）居住环境变量的影响}

城市与居住区域两组虚拟变量的系数\textbf{均不显著}（$P>0.05$），说明在控制噪声感知因子后，\textbf{行政区划与城乡类型本身并未带来独立的满意度效应}。

与之形成鲜明对比的是住宅的物理隔音条件。配备隔音墙使满意度处于更高等级的累积优势提升至无隔音墙居民的\textbf{1.953 倍}（$\beta=0.670$，$P<0.001$，95\%置信区间1.318-2.894），是居住环境类变量中\textbf{唯一显著}的一项。

这说明决定居民满意度的是\textbf{住宅自身的隔声性能}，而非宏观区位；噪声治理的着力点应落在建筑围护结构的改善上。

\textbf{（三）噪声感知因子的影响}

四个噪声感知公因子的系数\textbf{全部为负且均达到显著水平}，验证了“噪声影响感知越强、声环境满意度越低”这一核心假设。四项系数绝对值依次为$0.583$、$0.440$、$0.309$、$0.256$，呈现清晰的\textbf{梯度差异}。

影响力度最强的是 $F_2$ \textbf{情绪与社会关系影响因子}（$\beta=-0.583$，$P<0.001$）：因子得分每提高1个单位，满意度处于更高等级的累积优势降为原来的\textbf{55.8\%}。该维度直接关联居民对生活品质的主观评价，是决定满意度的\textbf{核心敏感维度}。

影响力度次之的是 $F_3$ \textbf{典型环境噪声源干扰因子}（$\beta=-0.440$，$P<0.001$，累积优势\textbf{64.4\%}），其后为 $F_4$ \textbf{生活场景噪声源干扰因子}（$\beta=-0.309$，$P=0.001$，累积优势\textbf{73.4\%}）。两者分别对应公共性噪声源与贴近个体生活的分散噪声源，负向作用明确，但强度均弱于情绪维度。

影响力度最弱的是 $F_1$ \textbf{身心健康与认知干扰因子}（$\beta=-0.256$，$P=0.005$，累积优势\textbf{77.4\%}）。该系数仍显著为负，说明在同时纳入情绪维度后，噪声对身心健康的干扰所对应的满意度差异相对有限。

\textbf{（四）影响因素的综合比较}

三类变量对满意度的解释力存在明显的\textbf{主次之分}：噪声感知因子是核心驱动变量，人口特征与居住环境变量的作用相对有限。

居住环境类变量中唯一稳定的是\textbf{物理隔音条件}，隔音墙的 $OR$ 值达到\textbf{1.953}，影响强度居全部非噪声感知变量之首。

城市与居住区域不显著这一阴性结果值得关注：三市满意度的原始分布虽有差异，但在控制噪声感知与隔音条件后差异消失，说明\textbf{城市间的差距源于各市居民噪声感知水平与住宅隔音条件的构成差异}，而非城市本身的独立效应。

人口特征变量中，\textbf{青少年（18岁及以下）}满意度最低，其次是46-60岁中年与大专学历群体，应作为差异化关注对象。

提升满意度的着力点可归纳为两个方向：\textbf{优先降低噪声对居民情绪与社会关系的干扰}，以及\textbf{推进隔音条件较差住宅的隔音改造}。

上述判断与聚类分析识别出的\textbf{“身心情绪高敏感型”重点人群}相互印证，构成本文后续治理建议的依据。

\subsection{模型稳健性检验}

为验证有序Logistic回归作为主要分析模型的合理性，本文将其与随机森林、XGBoost两种机器学习模型在独立的留出测试集（20\%样本）上进行对比。三种模型均采用分类形式，其中有序Logistic为有序分类，随机森林与XGBoost为多分类模型；XGBoost经贝叶斯调参（Optuna的TPE采样器，以宏平均F1值为优化目标，最优值0.3003）。对比结果如表\ref{tab:comparison}所示。

\begin{table}[H]
\centering
\caption{三种模型在独立测试集上的分类性能对比}
\label{tab:comparison}
\rowcolors{2}{tablight}{white}
\begin{tabular}{M{4cm}M{3cm}M{3cm}M{3cm}}
\bluetoprule
\rowcolor{tabhead}
\thead{模型} & \thead{准确率} & \thead{宏平均F1值} & \thead{Kappa系数} \\
\hline
有序Logistic & 0.447 & 0.374 & 0.218 \\
随机森林 & 0.436 & 0.360 & 0.200 \\
XGBoost（调参后） & 0.426 & 0.331 & 0.192 \\
\bluebottomrule
\end{tabular}
\end{table}

\begin{figure}[H]
\centering
\begin{subfigure}[b]{0.32\textwidth}
  \centering
  \includegraphics[width=\textwidth]{compare_confusion_logistic.png}
  \caption{有序Logistic回归}
\end{subfigure}
\hfill
\begin{subfigure}[b]{0.32\textwidth}
  \centering
  \includegraphics[width=\textwidth]{compare_confusion_randomforest.png}
  \caption{随机森林}
\end{subfigure}
\hfill
\begin{subfigure}[b]{0.32\textwidth}
  \centering
  \includegraphics[width=\textwidth]{compare_confusion_xgboost.png}
  \caption{XGBoost}
\end{subfigure}
\caption{三种模型的混淆矩阵}
\label{fig:cm}
\end{figure}

由表\ref{tab:comparison}可知，有序Logistic在准确率、宏平均F1值与Kappa三项指标上均优于随机森林与XGBoost，综合表现最优。图\ref{fig:cm}给出了三种模型在测试集上的混淆矩阵，三模型均呈现“对角线占优、误差集中于相邻等级”的形态，跨级误判较少，说明模型较好地保留了满意度的有序结构；其中有序Logistic的对角线集中程度最高，与其准确率、Kappa最高相一致。这表明在小样本有序分类场景下，可解释的有序Logistic模型并未因线性假设而损失预测性能，反而优于集成树模型，印证了本文将其作为主要分析模型的合理性。机器学习模型在本研究中的作用在于为主模型提供稳健性佐证，而非替代。

